\chapter{Практикум і задачі}\label{ch:geo7}

\section{Контрольні запитання}\label{sec:geo7-q}

\begin{enumerate}
\item Чому вся геометрія осесиметричної рівноваги задається однією функцією $\psi(R,Z)$? Що саме порушує осесиметрію в реальній машині?
\item Чому в таблицях і графіках поверхні позначають через $\psiN$, а не через $\psi$? Що станеться з $\psiN$, якщо помножити $\psi$ на $-1$?
\item Чому $\qnf$, а не $q$ на межі? Що відбувається з $q$ на сепаратрисі?
\item Чим кут $\thetastar$ кращий за геометричний $\theta$ для опису островів? Чому лінії сталого $\thetastar$ густіші на внутрішньому боці тора?
\item Чому $\thetastar$ не залежить від $F$, а $q$ залежить лінійно?
\item Чим відрізняються D Міллера (плазма) і D чистого натягу (котушка TF)? Що визначає кожну з них?
\item Чому гілки сепаратриси перетинаються майже під прямим кутом? Коли кут відрізнявся б від $90^\circ$?
\item Чому маска «$\psiN<1$» не збігається з областю «всередині LCFS» у діверторній плазмі?
\end{enumerate}

\section{Задачі}\label{sec:geo7-p}

\begin{enumerate}
\item\label{pr:geo7-miller} \emph{Міллер.} Доведіть, що в параметризації~\eqref{eq:geo4-miller} верхня точка має $R=R_0-a\delta$ і $Z=\kappa a$, тобто $\kappa$ і $\delta$ збігаються з означеннями~\eqref{eq:geo4-shape}. Обчисліть $R_{Z_{\max}}$ для JET ($R_0=\SI{2,96}{м}$, $a=\SI{1,25}{м}$, $\delta=0{,}3476$).
\item\label{pr:geo7-qcyl} \emph{Виграш від витягнутості.} За~\eqref{eq:geo4-qcyl} оцініть $q_{\mathrm{cyl}}$ для DIII-D \#203702 ($a=\SI{0,605}{м}$, $\kappa=1{,}70$, $R_{\mathrm{geo}}=\SI{1,656}{м}$, $\Btor(R_{\mathrm{geo}})=\SI{1,953}{Тл}$, $|\Ip|=\SI{1,359}{МА}$) з урахуванням і без урахування $\kappa$. Порівняйте з $\qnf=3{,}70$. Що може пояснити різницю?
\item\label{pr:geo7-pappus} \emph{Площа й об'єм.} Для JET порівняйте площу й об'єм межі Міллера (табл.~\ref{tab:geo4-shape}) з еліпсом тих самих $a$, $\kappa$ з центром на $R_0$. Чому трикутність зменшує об'єм помітніше, ніж площу?
\item\label{pr:geo7-axis} \emph{Вісь.} За гесіаном на осі DIII-D ($\psi_{RR}=3{,}473$, $\psi_{ZZ}=2{,}288$, $\psi_{RZ}=-0{,}036$~Вб/(рад$\cdot$м$^2$)), $F=\SI{3,2334}{Т.м}$ і $\Rax=\SI{1,6782}{м}$ знайдіть $\kappa_0$ і $q_0$ за~\eqref{eq:geo5-q0}. Виведіть~\eqref{eq:geo5-q0}.
\item\label{pr:geo7-x} \emph{X-точка.} Власні значення гесіана в X-точці DIII-D --- $-0{,}743$ і $+0{,}710$~Вб/(рад$\cdot$м$^2$), $R_X=\SI{1,241}{м}$. Знайдіть кут між гілками сепаратриси і густину струму $|\Jtor|$ в X-точці. Порівняйте з середньою $|\Ip|/S$.
\item\label{pr:geo7-coil} \emph{D-котушка.} Для $R_1=\SI{1,304}{м}$, $R_2=\SI{4,790}{м}$ знайдіть $\kD$ і радіуси кривини на кінцях дуги. Наскільки менший радіус біля внутрішньої ноги, ніж на зовнішньому екваторі?
\item\label{pr:geo7-ripple} \emph{Гофрування.} Оцініть $\dripple$ на $R=\SI{4,21}{м}$ і на $R_0=\SI{2,96}{м}$ як $(R/R_2)^N$ при $N=32$. Порівняйте з Біо--Саваром ($1{,}83\,\%$ і $3\cdot10^{-7}$). Чому на осі проста оцінка непридатна?
\end{enumerate}

\begin{practicum}[міні-проєкт]
Напишіть функцію, яка для будь-якого кадру DIII-D чи MAST за 1~с повертає «паспорт геометрії»: $\Rax$, $\Zax$, $R_{\mathrm{geo}}$, $a$, $A$, $\kappa$, $\dup$, $\dlo$, $S$, $V$, $|\Ip|$, координати й $\psiN$ X-точок, $q_0$ (з гесіана), $\qnf$. Усі цеглинки вже є в \texttt{geom\_tour.py} і \texttt{tokviz.equilibrium}. Перевірте на кадрі 75: паспорт має відтворити числа розд.~\ref{ch:geo2}--\ref{ch:geo5}.
\end{practicum}
